\documentclass[a4,10pt]{book}
%-----------------------------------------------------------
\usepackage[top=0.1in, bottom=0in, left=0.25in, right=0.25in]{geometry}
\usepackage{graphicx}
\usepackage{xcolor}
\usepackage{url}
\usepackage{palatino}
\usepackage{pdfrender}
\usepackage{tabularx}
\definecolor{darkish-blue}{rgb}{0,0,0.4}
\usepackage[colorlinks=true,linkcolor=black,anchorcolor=black,citecolor=black,filecolor=black,menucolor=black,runcolor=black,urlcolor=darkish-blue]{hyperref}
\usepackage{wrapfig}
\fontfamily{SansSerif}
\selectfont

\usepackage[T1]{fontenc}
\usepackage
%[ansinew]
[utf8]
{inputenc}

\usepackage{color}
\definecolor{mygrey}{gray}{0.75}
\textheight=9.75in
\raggedbottom

\setlength{\tabcolsep}{0in}
\newcommand{\isep}{-2 pt}
\newcommand{\lsep}{-0.57cm}
\newcommand{\psep}{-0.6cm}
\renewcommand{\labelitemii}{$\circ$}

\pagestyle{empty}
%-----------------------------------------------------------
%Custom commands
\newcommand{\resitem}[1]{\item #1 \vspace{-2pt}}
\newcommand{\resheading}[1]{{\small \colorbox{mygrey}{\begin{minipage}{0.975\textwidth}{\textbf{#1 \vphantom{p\^{E}}}}\end{minipage}}}}
\newcommand{\ressubheading}[3]{
\begin{tabular*}{6.62in}{l @{\extracolsep{\fill}} r}
	\textsc{{\textbf{#1}}} & \textsc{\textit{[#2]}} \\
\end{tabular*}\vspace{-8pt}}
%-----------------------------------------------------------
\begin{document}
\begin{center}
\textsc{\huge Aditya Rastogi}
\end{center}

\begin{minipage}{0.30\textwidth}
{\raggedleft
\href{https://github.com/thunderInfy/}{GitHub Profile} \\
\href{http://cse.iitkgp.ac.in/~arastogi}{Personal Website CSE - IITKGP}\\
}
\end{minipage}%
\begin{minipage}{0.33\textwidth}
\begin{center}
	\includegraphics[scale=0.1]{logo.jpg}
\end{center}
\end{minipage}%
\begin{minipage}{0.33\textwidth}
{\raggedright
Email-id : \textbf{r.aditya0824@gmail.com} \\
Mobile No.: \textbf{917550916164} \\
}
\end{minipage}%
\hspace{0.5cm}\\
\\
\resheading{\textbf{ACADEMIC DETAILS} }\\[\lsep]
\\
\begin{center}
\renewcommand{\arraystretch}{1.5}
\indent \begin{tabular}{ |@{\hskip 0.125in}l @{\hskip 0.125in} |@{\hskip 0.125in}l @{\hskip 0.125in} |@{\hskip 0.125in}l @{\hskip 0.125in} |@{\hskip 0.125in}l @{\hskip 0.25in} |l }
\hline
\textbf{Education} & \textbf{Institute} & \textbf{Year} & \textbf{CGPA / \%} \\
\hline
B. Tech and M. Tech (Dual Degree): &&&\\[-0.5em]
{Computer Science and Engineering} & IIT Kharagpur  & 2016 - 2021 (Expected) & \textbf{9.43 / 10} \\
\hline
\end{tabular}
\end{center}
\hfill 
\hspace{0.5cm}\\
\resheading{\textbf{INTERNSHIPS} }\\[\lsep]
\\\\
\large \href{http://netsyslab.ece.ubc.ca/wiki/index.php/People}{\textbf{Pattern Matching in Trillion Edge Graphs}} \normalsize \\
\emph{Mitacs Globalink Research Intern, University of British Columbia \hfill May'19-Jul'19 } \\ 
\emph{(Advisor: Prof. Matei Ripeanu)} \\[\lsep]
\begin{itemize}
\item Worked in a team to solve the problem of \textbf{approximate pattern matching in graphs} in a distributed systems setting.\\[\psep]
\item  Compared our pattern matching algorithm with TriAD, a distributed RDF engine.
\\[\lsep]
\end{itemize}
\emph{Research Areas:}  \textbf{Algorithms}, \textbf{Distributed Systems},
\textbf{Graph Theory}
\\\\
\large \textbf{Facial Landmarks Detection} \normalsize \\
\emph{Visiting Student Researcher, University of Sydney \hfill Dec'18-Jan'19} \\ 
\emph{(Advisor: Dr. Mehar Khatkar)} \\[\lsep]
\begin{itemize}
	\item  Designed and compared different \textbf{CNN} architectures with target point coordinates for detecting facial landmarks in different fish species.
	\\[\lsep]
\end{itemize}
\emph{Research Areas:}  \textbf{Deep Learning}, \textbf{Landmark Detection},
\textbf{Data Augmentation}
\\\\
\large \textbf{Sensor Diagnostics} \normalsize \\
\emph{Summer Intern, IIT-Kharagpur - Sponsored by Shell India Pvt. Ltd. \hfill May'18-July'18} \\ 
\emph{(Advisor: Prof. Swanand Khare)} \\[\lsep]
\begin{itemize}
\item  Developed a \textbf{Python application} to \textbf{sample from an arbitrary finite-ranged probability density} in one dimension.\\[\psep]
\item Reduced dimensionality of multiple sensors time-series data using \textbf{PCA} and \textbf{autoencoders}.\\[\psep]
\item Worked on \textbf{gaussian-mixture models} and error distributions in general.\\[\lsep]
\end{itemize}
\emph{Research Areas:}  \textbf{Machine Learning}, \textbf{Dimensionality Reduction},  \textbf{Predictive}  \textbf{Analysis}
\\
\hspace{0cm}\\[-0.5cm]\\
\resheading{\textbf{ PROJECTS } }\\[\lsep]
\\\\ \large \textbf{Deep Learning Visualization} \normalsize\\
\emph{TensorFlow.js, p5.js \hfill Aug'19 - Present} \\ 
\emph{(Advisors: Prof. Partha Pratim Chakrabarti and Prof. Aritra Hazra)} \\[\lsep]
\begin{itemize}
	\item Working on developing a meta algorithm which would visualize, understand and observe what neural networks learn in their underlying latent representations and takes decisions to alter the network to improve performance.\\[\psep]
	\item Working on actionability in neural networks given \textbf{feature visualizations} and \textbf{saliency mappings}.\\[\psep]
	\item Working on \textbf{adversarial robustness}, \textbf{calibration}, \textbf{out-of-distribution detection} and \textbf{energy based models}.
	\\
\end{itemize}
\large \href{http://cse.iitkgp.ac.in/~arastogi}{\textbf{Smart Rockets, Algorithmic Visualisations and Web-game Development}} \normalsize\\
\emph{p5-js, HTML, CSS} \\[\lsep]
\begin{itemize}
	\item Developed Smart Rockets, a demo of a \textbf{genetic algorithm}. Developed algorithmic visualisations on the web for algorithms like \textbf{Convex Hull}, Voronoi Diagram, Longest Increasing Subsequence etc.
	\\[\psep]
	\item  Used p5.js for creating \textbf{fractal} designs like trees, koch curves and flow field visualisation of \textbf{perlin noise}.\\[\psep]
	\item  Developed casual, fast-paced and 2D puzzle web-games.\\[\lsep]
\end{itemize}

\hspace{0.5cm}\\[-0.2cm]
\resheading{\textbf{BLOG POSTS} }\\[\lsep]
\\[-0.5em]
\begin{itemize}	
	\item 
	\href{https://medium.com/analytics-vidhya/understanding-simclr-a-simple-framework-for-contrastive-learning-of-visual-representations-d544a9003f3c?source=friends_link&sk=e768e4c2a88bd7d33f8eb42072b1fb05}{Understanding\textbf{ SimCLR} - A Simple Framework for \textbf{Contrastive Learning} of Visual Representations}
	\\[\lsep]
	\item 
	\href{https://towardsdatascience.com/solving-racetrack-in-reinforcement-learning-using-monte-carlo-control-bdee2aa4f04e?source=friends_link&sk=926f05a3660c6bf4d8c6a81be77645c5}{Solving Racetrack in \textbf{Reinforcement Learning} using \textbf{Monte Carlo Control}}
	\\[\lsep]
	\item
	\href{https://towardsdatascience.com/elucidating-policy-iteration-in-reinforcement-learning-jacks-car-rental-problem-d41b34c8aec7?source=friends_link&sk=6d58e79a8c87eb8396724901306312ac}{Elucidating \textbf{Policy Iteration} in \textbf{Reinforcement Learning} — Jack’s Car Rental Problem}
	\\[\lsep]
	\item 
	\href{https://medium.com/datadriveninvestor/visualizing-neural-networks-using-saliency-maps-in-pytorch-289d8e244ab4?source=friends_link&sk=37de912972716f67b3acab5c9bcdd9b0}{Visualizing Neural Networks using \textbf{Saliency Maps} in PyTorch}
	\\[\lsep]
	\item 
	\href{https://medium.com/swlh/the-gnu-toolchain-explained-4bf14666bc03?source=friends_link&sk=4b9c4214e37b2b38d98e8f79edc824fe}{The \textbf{GNU Toolchain} Explained}
\end{itemize}

\end{document}
